\chapter{Michal--Bastiani differentiability}
\label{chap:differentiability}

On a Banach space a map is $\Ck{1}$ when its derivative is a continuous map into the space of
bounded operators, with the operator norm. On a general locally convex space $E$ there is no
canonical topology on the space of continuous linear maps $E \to F$ that would play this role,
and the definition has to be rephrased. The solution of Michal \cite{Mic38} and Bastiani
\cite{Bas64} is to consider the derivative as a function of two variables, the point and the
direction, and to ask for joint continuity. In the recursive form used by Glöckner, which we
adopt, $\Ck{k+1}$ means that this function of two variables is $\Ck{k}$.

In this chapter we define $\Ck{k}$ maps between locally convex spaces and prove the basic
results of first-order calculus. After the definitions, we prove the mean value theorem, which
in this setting takes the form of an inclusion in a closed convex set. From it we obtain
Glöckner's lemma on the continuous extension of difference quotients, and from this lemma the
linearity of the differential and the chain rule. We close the chapter with partial
derivatives, the symmetry of the second differential, and the comparison with Fréchet $\Ck{k}$
maps on normed spaces. The references are \cite[\S1.2]{GN26} and \cite[\S I.2]{Nee06}.

\section{Directional derivatives and $\Ck{k}$ maps}

\begin{definition}[Directional derivative]
  \label{def:dir-deriv}
  \lean{LocallyConvex.Analysis.HasDirDerivAt, LocallyConvex.Analysis.dMB}
  \leanok
  Let $f \colon E \to F$ and $x, h \in E$. We say that $f$ has derivative $y \in F$ at $x$ in
  the direction $h$ if
  \[
    \lim_{t \to 0,\ t \neq 0} \frac{1}{t}\bigl(f(x + t h) - f(x)\bigr) = y .
  \]
  We write $\dd f(x)(h)$ for this limit when it exists, and regard $\dd f$ as a function on
  $E \times E$.
\end{definition}

The limit is taken in $F$, which is Hausdorff, so it is unique when it exists. Note that
$\dd f(x)(0) = 0$ always, and that no linearity in $h$ is built into the definition.

\begin{remark}
  In Lean, $\dd f$ is a total function: it takes an arbitrary value where the limit does not
  exist. This is harmless, since it is only used where the limit is known to exist.
\end{remark}

\begin{definition}[$\Ck{k}$ and smooth maps]
  \label{def:ck}
  \lean{LocallyConvex.Analysis.ContDiffMBOn, LocallyConvex.Analysis.SmoothMBOn}
  \leanok
  \uses{def:dir-deriv}
  Let $U \subseteq E$ be open. A map $f$ is $\Ck{0}$ on $U$ if it is continuous on $U$. It is
  $\Ck{k+1}$ on $U$ if it is continuous on $U$, all directional derivatives $\dd f(x)(h)$ exist
  for $x \in U$, $h \in E$, and $\dd f \colon U \times E \to F$ is $\Ck{k}$ on the open set
  $U \times E$. The map $f$ is smooth on $U$ if it is $\Ck{k}$ for every $k$.
\end{definition}

Thus $f$ is $\Ck{1}$ when its directional derivatives exist and are jointly continuous in the
point and the direction. This is \cite[1.2.1 and 1.3.1]{GN26}, in the recursive form used by Glöckner;
Bastiani's original definition, through the continuity of all the iterated differentials, is
equivalent (Theorem~\ref{thm:ck-iff-bastiani}).

\begin{remark}
  Two features of the definition deserve a comment. First, linearity of $h \mapsto \dd f(x)(h)$
  is not assumed; it is a theorem (Theorem~\ref{thm:linear}), and it depends on the joint
  continuity in an essential way (Example~\ref{ex:not-linear} below). Second, no topology on
  the space of continuous linear maps $E \to F$ is involved: the regularity of $\dd f$ is
  measured on the locally convex space $E \times E$, so that the induction stays within the
  category of locally convex spaces.
\end{remark}

\begin{remark}
  The recursive definition quantifies over the domain of $\dd f$, which is $U \times E$ rather
  than $U$. As a consequence, in Lean the domains of $\Ck{k}$ maps live in a fixed universe.
  Codomains are unrestricted.
\end{remark}

We record the two basic examples, which are used throughout.

\begin{lemma}[Continuous linear maps are smooth]
  \label{lem:clm-smooth}
  \lean{ContinuousLinearMap.contDiffMBOn, ContinuousLinearMap.smoothMBOn}
  \leanok
  \uses{def:ck}
  Every continuous linear map $L \colon E \to F$ is smooth, and $\dd L(x)(h) = L h$.
\end{lemma}

\begin{proof}
  \leanok
  The difference quotients of $L$ are constant, equal to $Lh$. Then $\dd L = L \circ \pi_2$ is
  again continuous and linear, and the claim follows by induction on $k$.
\end{proof}

\begin{lemma}[Pairs]
  \label{lem:prod}
  \lean{LocallyConvex.Analysis.ContDiffMBOn.prodMk}
  \leanok
  \uses{def:ck}
  If $f \colon U \to F$ and $g \colon U \to G$ are $\Ck{k}$, so is $x \mapsto (f(x), g(x))$.
\end{lemma}

\begin{proof}
  \leanok
  Directional derivatives of a pair are computed componentwise, so $\dd (f,g)$ is the pair
  $(\dd f, \dd g)$; the claim follows by induction on $k$.
\end{proof}

\section{The mean value theorem}

The classical mean value inequality $\|f(x+h) - f(x)\| \leq \sup_s \|Df(x+sh)\|\,\|h\|$ has no
meaning without a norm. Its replacement is an inclusion: increments of $f$ lie in the closed
convex hull of the derivatives along the segment. Since closed convex neighborhoods form a
basis of the topology, this inclusion carries the same information as a family of estimates,
one for each continuous seminorm.

\begin{lemma}[Mean value inclusion]
  \label{lem:mean-value-inclusion}
  \lean{LocallyConvex.Analysis.sub_mem_of_hasDirDerivAt}
  \leanok
  \uses{def:dir-deriv}
  Let $z, w \in E$ and suppose that $f$ has derivative $D(s)$ at $z + s w$ in the direction $w$
  for every $s \in [0,1]$, with $D$ continuous on $[0,1]$. If $W \subseteq F$ is closed, convex
  and contains $D(s)$ for every $s \in [0,1]$, then $f(z + w) - f(z) \in W$.
\end{lemma}

\begin{proof}
  \leanok
  Suppose not. By Hahn--Banach there are a continuous linear functional $\varphi$ and $c \in \R$
  with $\varphi(f(z+w) - f(z)) > c \geq \varphi(v)$ for all $v \in W$. The real function
  $g(s) = \varphi(f(z + s w))$ has derivative $\varphi(D(s)) \leq c$ on $[0,1]$, and this
  derivative is continuous. By the fundamental theorem of calculus for real functions,
  $g(1) - g(0) = \int_0^1 \varphi(D(s))\,ds \leq c$, a contradiction.
\end{proof}

This is the remark after Lemma~8.4 in \cite{Mic80}. A first consequence is that the continuity
required in Definition~\ref{def:ck} is, from the first order on, automatic.

\begin{lemma}[Continuity]
  \label{lem:continuity}
  \lean{LocallyConvex.Analysis.continuousOn_of_hasDirDerivAt}
  \leanok
  \uses{lem:mean-value-inclusion}
  If all directional derivatives of $f$ exist on $U$ and are jointly continuous on $U \times E$,
  then $f$ is continuous on $U$.
\end{lemma}

\begin{proof}
  \leanok
  Since $\dd f(x)(0) = 0$, joint continuity gives, for each closed convex neighborhood $W$ of
  $0$, a neighborhood of $(x, 0)$ on which $\dd f$ takes values in $W$. For $w$ small, the
  mean value inclusion then gives $f(x + w) - f(x) \in W$. This is \cite[8.1]{Mic80}.
\end{proof}

When the derivatives are tested against continuous linear functionals, the mean value inclusion
becomes an identity.

\begin{proposition}[Mean value theorem, weak form]
  \label{prop:mean-value-weak}
  \lean{LocallyConvex.Analysis.ContDiffMBOn.integral_dMB}
  \leanok
  \uses{def:ck}
  Let $f$ be $\Ck{1}$ on a set containing the segment $\{x + t h : t \in [0,1]\}$. Then, for
  every continuous linear functional $\varphi$ on $F$,
  \[
    \varphi\bigl(f(x+h)\bigr) - \varphi\bigl(f(x)\bigr)
      = \int_0^1 \varphi\bigl(\dd f(x + t h)(h)\bigr)\, dt .
  \]
\end{proposition}

\begin{proof}
  \leanok
  The real function $t \mapsto \varphi(f(x + t h))$ has continuous derivative
  $t \mapsto \varphi(\dd f(x + th)(h))$; it suffices to apply the fundamental theorem of
  calculus for real functions.
\end{proof}

In the language of Chapter~\ref{chap:integrals}, $f(x+h) - f(x)$ is the weak integral of
$t \mapsto \dd f(x + th)(h)$ over $[0,1]$ (Corollary~\ref{cor:ftc-segment}). This is \cite[1.2.6]{GN26}.
Note that no completeness of $F$ is needed here: the integral exists because its value,
$f(x+h) - f(x)$, is given in advance.

\section{Glöckner's lemma and the linearity of the differential}

The following lemma is the technical core of the chapter. It says that, for a $\Ck{1}$ map, the
difference quotient is jointly continuous in the base point, the direction and the step, up to
and including $t = 0$.

\begin{lemma}[Glöckner]
  \label{lem:glockner}
  \lean{LocallyConvex.Analysis.ContDiffMBOn.tendsto_quotient}
  \leanok
  \uses{def:ck, lem:mean-value-inclusion}
  Let $g$ be $\Ck{1}$ on $U$ and $y \in U$, $v \in E$. Then
  \[
    \frac{1}{t}\bigl(g(y' + t v') - g(y')\bigr) \longrightarrow \dd g(y)(v)
    \qquad \text{as } (y', v', t) \to (y, v, 0),\ t \neq 0 .
  \]
\end{lemma}

\begin{proof}
  \leanok
  By the mean value inclusion, the difference quotient lies in the closed convex hull of
  $\{\dd g(y' + s t v')(v') : s \in [0,1]\}$. Let $W$ be a closed convex neighborhood of
  $\dd g(y)(v)$. By continuity of $\dd g$ at $(y, v)$, all these derivatives lie in $W$ when
  $(y', v', t)$ is close to $(y, v, 0)$; hence so does the difference quotient. Since closed
  convex neighborhoods form a basis, the claim follows.
\end{proof}

In other words, the difference quotient extends continuously to $t = 0$ by $\dd g$. This is the
characterization of $\Ck{1}$ maps of Glöckner \cite{Glo02} (\cite[1.2.10]{GN26}), and it is the key
step in the chain rule. It also gives the linearity of the differential.

\begin{theorem}[The differential is linear]
  \label{thm:linear}
  \lean{LocallyConvex.Analysis.ContDiffMBOn.exists_clm}
  \leanok
  \uses{lem:glockner}
  If $f$ is $\Ck{1}$ on $U$ and $x \in U$, then $h \mapsto \dd f(x)(h)$ is a continuous linear
  map $E \to F$.
\end{theorem}

\begin{proof}
  \leanok
  Homogeneity, $\dd f(x)(sh) = s\, \dd f(x)(h)$, follows from the definition by a change of
  variable in the limit. For additivity we write
  \[
    \tfrac{1}{t}\bigl(f(x + th + tk) - f(x)\bigr)
    = \tfrac{1}{t}\bigl(f((x + th) + tk) - f(x + th)\bigr)
      + \tfrac{1}{t}\bigl(f(x + th) - f(x)\bigr) .
  \]
  As $t \to 0$ the second term tends to $\dd f(x)(h)$, and by Glöckner's lemma, applied at the
  moving base point $x + th \to x$, the first one tends to $\dd f(x)(k)$. Continuity follows from
  the joint continuity of $\dd f$. This is \cite[1.2.11]{GN26}.
\end{proof}

\begin{example}
  \label{ex:not-linear}
  The existence of all directional derivatives, even everywhere and for a continuous map, does
  not imply linearity. Let $f \colon \R^2 \to \R$ be given by $f(a, b) = a^2 b/(a^2 + b^2)$ for
  $(a,b) \neq 0$ and $f(0) = 0$. It is continuous, since $|f(a,b)| \leq |b|$, and smooth away
  from the origin. Since $f(t a, t b) = t f(a, b)$, we have $\dd f(0)(a, b) = f(a, b)$, which
  is not additive: $\dd f(0)(1,0) = \dd f(0)(0,1) = 0$, while $\dd f(0)(1,1) = 1/2$. By
  Theorem~\ref{thm:linear}, $f$ is not $\Ck{1}$; since $f$ is smooth away from the origin,
  what fails is the joint continuity of $\dd f$ at some point of the form $(0, h)$.
\end{example}

\section{The chain rule}

\begin{theorem}[Chain rule]
  \label{thm:chain-rule}
  \lean{LocallyConvex.Analysis.ContDiffMBOn.comp, LocallyConvex.Analysis.ContDiffMBOn.dMB_comp}
  \leanok
  \uses{lem:glockner, lem:prod}
  Let $U \subseteq E$ and $V \subseteq F$ be open, $f \colon U \to F$ with $f(U) \subseteq V$
  and $g \colon V \to G$. If $f$ and $g$ are $\Ck{k}$, then $g \circ f$ is $\Ck{k}$, and for
  $k \geq 1$
  \[
    \dd (g \circ f)(x)(h) = \dd g\bigl(f(x)\bigr)\bigl(\dd f(x)(h)\bigr) .
  \]
\end{theorem}

\begin{proof}
  \leanok
  For $t \neq 0$ write $f(x + th) = f(x) + t\, q(t)$, where
  $q(t) = \frac{1}{t}(f(x + th) - f(x)) \to \dd f(x)(h)$. Then
  \[
    \tfrac{1}{t}\bigl(g(f(x + th)) - g(f(x))\bigr)
    = \tfrac{1}{t}\bigl(g(f(x) + t\, q(t)) - g(f(x))\bigr),
  \]
  which tends to $\dd g(f(x))(\dd f(x)(h))$ by Glöckner's lemma at
  $(y', v', t) = (f(x), q(t), t)$. This gives the formula. It says that
  $\dd (g \circ f) = \dd g \circ (f \circ \pi_1, \dd f)$, a composition of $\Ck{k-1}$ maps, so
  the regularity follows by induction on $k$. This is \cite[1.2.13 and 1.3.4]{GN26}.
\end{proof}

\begin{remark}
  The direction $q(t)$ at which Glöckner's lemma is applied moves with $t$. This is where the
  continuity of $\dd g$ in the direction, and not only in the base point, enters the proof.
\end{remark}

\section{Partial derivatives}

As in finite dimension, a map on a product is $\Ck{1}$ as soon as its partial derivatives exist
and are continuous.

\begin{theorem}[Continuous partial derivatives]
  \label{thm:partials}
  \lean{LocallyConvex.Analysis.contDiffMBOn_one_of_partials}
  \leanok
  \uses{lem:mean-value-inclusion, def:ck}
  Let $U \subseteq E \times F$ be open and $f \colon U \to G$. Suppose that the partial
  derivatives $\dd_1 f(p)(u) = \dd f(p)(u, 0)$ and $\dd_2 f(p)(v) = \dd f(p)(0, v)$ exist for
  $p \in U$, and are jointly continuous on $U \times E$ and $U \times F$. Then $f$ is $\Ck{1}$ on
  $U$, and $\dd f(p)(u, v) = \dd_1 f(p)(u) + \dd_2 f(p)(v)$.
\end{theorem}

\begin{proof}
  \leanok
  We split
  \[
    f(p + t(u,v)) - f(p) = \bigl[f(p + t(u,v)) - f(p + t(0,v))\bigr]
      + \bigl[f(p + t(0,v)) - f(p)\bigr].
  \]
  Divided by $t$, the second bracket tends to $\dd_2 f(p)(v)$. By the mean value inclusion, the
  first one, divided by $t$, lies in the closed convex hull of
  $\{\dd_1 f(p + t(0,v) + st(u,0))(u) : s \in [0,1]\}$, and these values tend to
  $\dd_1 f(p)(u)$ by joint continuity. The resulting formula for $\dd f$ is jointly continuous,
  and $f$ is continuous by Lemma~\ref{lem:continuity}. This is \cite[8.3]{Mic80} and \cite[1.2.8]{GN26}.
\end{proof}

This criterion is a practical way to verify that a given map is $\Ck{1}$; for instance, it is
used in the proof that continuous bilinear maps are smooth (Theorem~\ref{thm:bilinear}).

\section{Symmetry of the second differential}

For a $\Ck{2}$ map, the second differential at $x$ is
$(h_1, h_2) \mapsto \dd(\dd f)\bigl((x, h_1), (h_2, 0)\bigr)$, the derivative of
$\dd f(\cdot)(h_1)$ in the direction $h_2$. We prove that it is symmetric. As in the classical
proof, both orders of differentiation are compared with a second difference quotient, which is
symmetric by construction.

\begin{lemma}[Second differences]
  \label{lem:second-diff}
  \lean{LocallyConvex.Analysis.ContDiffMBOn.second_diff_mem}
  \leanok
  \uses{lem:mean-value-inclusion, thm:linear}
  Let $f$ be $\Ck{2}$ on $U$, $x \in U$ and $h_1, h_2 \in E$. For every closed convex
  neighborhood $W$ of $\dd(\dd f)\bigl((x, h_1), (h_2, 0)\bigr)$, the second difference quotient
  \[
    \frac{1}{t^2}\bigl(f(x + t h_1 + t h_2) - f(x + t h_1) - f(x + t h_2) + f(x)\bigr)
  \]
  lies in $W$ for all small $t \neq 0$.
\end{lemma}

\begin{proof}
  \leanok
  We apply the mean value inclusion twice: first to $s \mapsto f(x + s h_1 + t h_2) - f(x + s h_1)$,
  then to the derivative of this map in the direction $h_2$. The second difference quotient
  lies in the closed convex hull of values of the second differential at points near $x$, and
  these lie in $W$ by continuity.
\end{proof}

\begin{theorem}[Schwarz, second order]
  \label{thm:symm2}
  \lean{LocallyConvex.Analysis.ContDiffMBOn.dMB_dMB_symm}
  \leanok
  \uses{lem:second-diff}
  If $f$ is $\Ck{2}$ on $U$ and $x \in U$, then
  $\dd(\dd f)\bigl((x, h_1), (h_2, 0)\bigr) = \dd(\dd f)\bigl((x, h_2), (h_1, 0)\bigr)$ for all
  $h_1, h_2 \in E$.
\end{theorem}

\begin{proof}
  \leanok
  The second difference quotient is symmetric in $h_1$ and $h_2$. By the previous lemma it
  converges to both sides, and limits in a Hausdorff space are unique. This is \cite[1.3.17]{GN26} for
  $k = 2$.
\end{proof}

The symmetry of the higher differentials is the subject of Chapter~\ref{chap:higher}.

\section{Fréchet differentiable maps}

On normed spaces both calculi are available, and the Fréchet one is the stronger of the two.

\begin{proposition}[Fréchet $\Ck{k}$ implies Michal--Bastiani $\Ck{k}$]
  \label{prop:frechet-to-mb}
  \lean{LocallyConvex.Analysis.ContDiffOn.contDiffMBOn, HasFDerivAt.hasDirDerivAt}
  \leanok
  \uses{def:ck}
  Let $E$ and $F$ be normed spaces. If $f$ is $\Ck{k}$ on $U$ in the sense of Fréchet, then it
  is $\Ck{k}$ in the sense of Michal--Bastiani.
\end{proposition}

\begin{proof}
  \leanok
  A Fréchet differentiable map has all directional derivatives, with $\dd f(x)(h) = Df(x)h$.
  The map $(x, h) \mapsto Df(x) h$ is Fréchet $\Ck{k-1}$ on $U \times E$, and the claim follows by
  induction on $k$. This is \cite[Remark I.2.2]{Nee06}.
\end{proof}

The converse fails for $\Ck{k}$, but holds up to the loss of one derivative; see
Chapter~\ref{chap:frechet}.
